\documentclass[11pt,a4paper]{article}
\usepackage[utf8]{inputenc}
\usepackage[T1]{fontenc}
\usepackage[portuguese]{babel}
\usepackage{geometry}
\geometry{margin=2.5cm}
\usepackage{listings}
\usepackage{xcolor}
\usepackage{amsmath}
\usepackage{amssymb}
\usepackage{hyperref}
\usepackage{enumitem}
\usepackage{tcolorbox}
\usepackage{tabularx}
\usepackage{booktabs}

\definecolor{codebg}{rgb}{0.95,0.95,0.95}
\definecolor{codegreen}{rgb}{0.1,0.5,0.1}
\definecolor{codepurple}{rgb}{0.5,0.1,0.5}
\definecolor{codegray}{rgb}{0.5,0.5,0.5}

\lstdefinestyle{python}{
    backgroundcolor=\color{codebg},
    commentstyle=\color{codegray}\itshape,
    keywordstyle=\color{codepurple}\bfseries,
    stringstyle=\color{codegreen},
    basicstyle=\ttfamily\small,
    breaklines=true,
    frame=single,
    language=Python,
    showstringspaces=false,
    tabsize=4,
    literate={á}{{\'a}}1 {ã}{{\~a}}1 {é}{{\'e}}1 {í}{{\'i}}1 {ó}{{\'o}}1 {õ}{{\~o}}1 {ú}{{\'u}}1 {ç}{{\c{c}}}1 {Á}{{\'A}}1 {Ã}{{\~A}}1 {É}{{\'E}}1 {Í}{{\'I}}1 {Ó}{{\'O}}1 {Õ}{{\~O}}1 {Ú}{{\'U}}1 {Ç}{{\c{C}}}1 {â}{{\^a}}1 {ê}{{\^e}}1 {ô}{{\^o}}1 {û}{{\^u}}1 {Â}{{\^A}}1 {Ê}{{\^E}}1 {Ô}{{\^O}}1 {Û}{{\^U}}1 {à}{{\`a}}1 {è}{{\`e}}1 {ì}{{\`i}}1 {ò}{{\`o}}1 {ù}{{\`u}}1
}
\lstset{style=python}

\newtcolorbox{objetivos}[1][]{
  colback=green!5!white,
  colframe=green!40!black,
  title=O que você vai aprender nesta página,
  fonttitle=\bfseries,
  #1
}

\newtcolorbox{exercicio}[1]{
  colback=blue!5!white,
  colframe=blue!40!black,
  title=#1,
  fonttitle=\bfseries
}

\newtcolorbox{solucao}{
  colback=yellow!5!white,
  colframe=orange!60!black,
  title=Solução proposta,
  fonttitle=\bfseries
}

\title{\textbf{Data Analysis with Python 2024--2025}\\Parte 2: NumPy (Parte 1)\\\large Caderno de Exercícios}
\author{Tradução e Adaptação: University of Helsinki --- MOOC.fi}
\date{}

\begin{document}

\maketitle

\section*{Nota Importante}
Este material é uma tradução e adaptação baseada no curso \textit{Data Analysis with Python 2024-2025}, oferecido pela \textbf{Universidade de Helsinque (MOOC.fi)}.

\tableofcontents
\newpage

\section{Exercícios - Capítulo 2 (NumPy)}

\subsection*{Exercício 2.11 -- Linhas e colunas}

\begin{exercicio}{Sumário}
\begin{itemize}
    \item \textbf{Objetivo:} Escrever as funções \texttt{get\_rows} e \texttt{get\_columns} que recebam uma matriz bidimensional e retornem, respectivamente, listas contendo as linhas e colunas formatadas como vetores unidimensionais.
    \item \textbf{Dica:} A operação de transposição (método \texttt{.T}) pode facilitar a extração das colunas ao ser combinada com a função de linhas.
\end{itemize}
\end{exercicio}

\subsection*{Exercício 2.12 -- Vetores-linha e vetores-coluna}

\begin{exercicio}{Sumário}
\begin{itemize}
    \item \textbf{Objetivo:} Extrair linhas e colunas de uma matriz 2D de formato \texttt{(n, m)}, mas mantendo estritamente a estrutura 2D original dos vetores.
    \item \textbf{Formato Esperado:} A função \texttt{get\_row\_vectors} deve retornar vetores no formato \texttt{(1, m)}. A função \texttt{get\_column\_vectors} deve retornar vetores no formato \texttt{(n, 1)}. Utilize \textit{slicing} inteligente que preserva a dimensão.
\end{itemize}
\end{exercicio}

\subsection*{Exercício 2.13 -- Diamante}

\begin{exercicio}{Sumário}
\begin{itemize}
    \item \textbf{Objetivo:} Criar a função \texttt{diamond(n)} que gera uma matriz bidimensional representando a figura de um losango/diamante delimitado por valores \texttt{1}, enquanto os outros valores permanecem \texttt{0}.
    \item \textbf{Ferramentas:} Utilize a manipulação de blocos de matrizes unindo \texttt{np.eye} (para a geração de diagonais), operações de reflexão (\texttt{np.fliplr}, \texttt{np.flipud}), e junção com \texttt{np.concatenate}.
\end{itemize}
\end{exercicio}

\subsection*{Exercício 2.14 -- Comprimento de vetores}

\begin{exercicio}{Sumário}
\begin{itemize}
    \item \textbf{Objetivo:} Criar a função \texttt{vector\_lengths(a)} que receba uma matriz de formato \texttt{(n, m)} e calcule o comprimento/norma euclidiana de cada uma das \texttt{n} linhas.
    \item \textbf{Restrições:} É estritamente \textbf{proibido o uso de laços} (\texttt{for}/\texttt{while}) e não é permitido invocar \texttt{scipy.linalg.norm}.
    \item \textbf{Ferramentas Matemáticas:} A solução deve ser vetorizada, baseando-se no fato que a norma Euclidiana é a raiz da soma dos quadrados. Utilize pelo menos as funções \texttt{np.sum} e \texttt{np.sqrt}.
\end{itemize}
\end{exercicio}

\subsection*{Exercício 2.15 -- Ângulos entre vetores}

\begin{exercicio}{Sumário}
\begin{itemize}
    \item \textbf{Objetivo:} Desenvolver \texttt{vector\_angles(X, Y)} para calcular simultaneamente os ângulos (em graus) entre vetores representados pelas linhas das matrizes \texttt{X} e \texttt{Y}.
    \item \textbf{Fórmula:} O cosseno do ângulo é obtido dividindo o produto interno escalar dos vetores pelo produto numérico das suas respectivas normas.
    \item \textbf{Restrições:} O uso de laços é proibido. Para contornar imperfeições e imprecisões no processamento em float, antes de passar o cosseno para o \texttt{np.arccos}, garanta que ele esteja preso no limite entre \texttt{[-1.0, 1.0]} utilizando o método \texttt{np.clip}. Aplique \texttt{np.degrees} na saída.
\end{itemize}
\end{exercicio}

\subsection*{Exercício 2.16 -- Tabuada revisitada}

\begin{exercicio}{Sumário}
\begin{itemize}
    \item \textbf{Objetivo:} Construir a função \texttt{multiplication\_table(n)} que retorne um array bidimensional de tamanho \texttt{(n, n)}, em que cada célula \texttt{i, j} possui o valor exato da multiplicação de seus eixos (\texttt{i * j}).
    \item \textbf{Restrições:} Não utilizar laços \texttt{for} na solução final.
    \item \textbf{Ferramentas:} Sua resolução obrigatoriamente exigirá a aplicação do conceito de \textit{broadcasting}. Instancie os multiplicadores base via \texttt{np.arange}, altere seus eixos utilizando \texttt{reshape}, e delegue a carga matemática aos operadores universais do NumPy.
\end{itemize}
\end{exercicio}

\end{document}
